% Clausura focal de antecedentes: fuentes del artículo XI, edición 21-09-2026.
% Rangos, copias íntegras y huellas: gestion/mapa_restitucion_memoria.json.
\section{Antécédents de restitution opératoire et de mémoire}
\label{md:antmem:capitulo}

Le point de départ de ces antécédents se situe après la production conjointe
APP--TRIT--TPK de l’état enrichi et de la structure discrète du
continu. On travaille sur le registre dodécaphasique déjà engendré et sur des
réalisations déclarées de ses lecteurs. Le support linéaire ultérieur
et les candidats d’une reconstruction ne remplacent pas le domaine des
histoires engendrées. Une réalisation réversible ultérieure n’est pas non plus
identifiée, par sa seule notation, à l’holonomie nonadique \(\Gamma_9\).

Dans la coordinatisation et l’orientation du registre,
\begin{equation}
 \begin{gathered}
 K=(234,543,140,729,659,824,621,58,914,794,146,601),\\
 \mathbf1^{*}K=6263,\qquad
 (Sx)_i=x_{i+1},\quad i\pmod {12}.
 \end{gathered}
 \label{md:antmem:registro}
\end{equation}
Le registre est utilisé après sa génération, non comme entrée
rétrospective pour sélectionner une sortie. L’opérateur \(S\) est ici
le décalage positionnel marqué. Ses domaines seront maintenus
séparés des plans canoniques utilisés plus loin.

\subsection{L’orbite complète comme référence de son porteur}
\label{md:antmem:k:orbita}

Dans \(V=\mathbb R^{12}\), muni de son produit euclidien canonique, considérons
\[
 O_K=(K,SK,\ldots,S^{11}K):\mathbb R^{12}\longrightarrow V.
\]
Le décalage \(S\) agit sur les positions du registre. Il n’est pas
identifié par la seule périodicité à l’évolution complète TPK ou à
une symétrie exceptionnelle ultérieure.

\begin{theorem}[Référence cyclique complète et matrice de Gram exacte]
\label{md:antmem:k:marco-ciclico}
La matrice \(O_K\) est inversible. Sa matrice de Gram a pour valeurs propres
\[
\begin{array}{c|c}
 \text{valeur}&\text{multiplicité}\\\hline
 39225169=6263^2&1\\
 1248025-345420\sqrt3&2\\
 589609&2\\
 1407529&2\\
 1053157&2\\
 1248025+345420\sqrt3&2\\
 697225&1
\end{array}
\]
et satisfait
\begin{equation}
 589609\|c\|^2\leq\|O_Kc\|^2\leq39225169\|c\|^2,
 \quad
 \|O_K^{-1}\|=\frac1{\sqrt{589609}},\quad
 \operatorname{cond}(O_K)=\frac{6263}{\sqrt{589609}}.
 \label{md:antmem:k:cotas-orbita}
\end{equation}
L’orbite du registre centré a rang onze.
\end{theorem}
\begin{proof}
Les entrées de la matrice de Gram sont les autocorrélations \(\langle K,S^jK\rangle\). Dans l'ordre \(j=0,\ldots,11\), elles valent
\begin{align*}
 (&4251257,2999323,3163385,3287920,3216553,3344743,\\
  &2950064,3344743,3216553,3287920,3163385,2999323).
\end{align*}
Comme \(S\) est orthogonal, le gramien est circulant. Substituer les douze
caractères \(\exp(2\pi i m/12)\) dans cette ligne donne le tableau réel
précédent. Toutes les valeurs sont positives. La plus petite des deux valeurs
avec \(\sqrt3\) est supérieure à \(589609\), car
\(658416>345420\sqrt3\), inégalité qui se vérifie en élevant au
carré ses deux membres positifs. Les autres valeurs du tableau sont
supérieures ou égales à \(589609\). Le mode uniforme a pour valeur
\((\sum K_i)^2=6263^2\) ; l’inégalité triangulaire de la transformée
de Fourier, en utilisant \(K_i\geq0\), borne toutes les autres par cette même
valeur. Le théorème spectral donne \eqref{md:antmem:k:cotas-orbita}. Centrer
le registre élimine seulement le mode uniforme ; les onze autres conservent
les valeurs non nulles du tableau.
\end{proof}

Le déterminant, avec l’ordre des colonnes et l’orientation de \(S\)
fixés ici, est
\[
 \det O_K=5483119259214020183850397930008283625.
\]
L'absence de période propre du vecteur ne prouverait pas à elle seule l'inversibilité : le point décisif est qu'aucun de ses modes de Fourier ne s'annule. Le conditionnement numérique est approximativement \(8.15643\). Si l'on normalise \(K\) à une norme égale à un, les extrêmes de la matrice de Gram sont divisés par \(\|K\|^2=4251257\) ; l'amplitude de ses composantes ne doit pas être confondue avec une amélioration intrinsèque par rapport à toute référence de même norme.

\begin{corollary}[Séparation polaire de la forme et de l’orientation]
\label{md:antmem:k:polar}
Soient
\[
 G_K=O_K^*O_K,\qquad P_K=G_K^{1/2},\qquad
 F_K=O_KG_K^{-1/2}.
\]
Alors \(P_K\) est défini positif, \(F_K\) est orthogonal et
\(O_K=F_KP_K\). Une réalisation isométrique \(J:V\to\mathcal Y\)
conserve \(G_K\) et \(P_K\), et transporte la partie isométrique vers
\(JF_K\).
\end{corollary}
\begin{proof}
La positivité du Gram permet de définir sa racine et son inverse au moyen du
théorème spectral. L’identité
\(F_K^*F_K=G_K^{-1/2}G_KG_K^{-1/2}=I\) prouve l’assertion.
De plus, \((JO_K)^*(JO_K)=O_K^*O_K\).
\end{proof}

Cette polaire est définie à partir du gramien complet : elle ne s’obtient pas seulement à partir de la
liste de ses valeurs propres ni de la norme de \(K\). Le symbole
\(P_K\) de ce corollaire n’est pas le projecteur dimensionnel de rang dix
d’une autre réalisation incidentielle. Cette séparation est une conséquence de la
référence cyclique, non une identification de tous ses lecteurs.

\subsection{Identification d'opérateurs et holonomie relative}
\label{md:antmem:k:identificacion}

\begin{theorem}[Douze réponses générales ou une réponse cyclique]
\label{md:antmem:k:operadores}
Pour tout opérateur linéaire \(H:V\to V\), ses douze réponses \(Y=(HK,HSK,\ldots,HS^{11}K)\) déterminent
\[
 H=YO_K^{-1}.
\]
Si \(HS=SH\), une unique réponse vectorielle complète \(y=HK\) détermine
\begin{equation}
 H=O_yO_K^{-1},\qquad
 H=\sum_{j=0}^{11}h_jS^j,\qquad h=O_K^{-1}y.
 \label{md:antmem:k:filtro}
\end{equation}
Avec des erreurs \(E\) dans la matrice des réponses ou \(\delta y\) dans le
vecteur, respectivement,
\begin{align}
 \|\widehat H-H\|&\leq\frac{\|E\|}{\sqrt{589609}},
 \label{md:antmem:k:error-matriz}\\
 \|\widehat h-h\|&\leq\frac{\|\delta y\|}{\sqrt{589609}},\qquad
 \|\widehat H-H\|\leq\sqrt{\frac{12}{589609}}\,\|\delta y\|.
 \label{md:antmem:k:error-filtro}
\end{align}
\end{theorem}
\begin{proof}
La première égalité est \(Y=HO_K\). Si \(H\) commute avec \(S\), alors \(HS^jK=S^jy\), de sorte que \(Y=O_y\). La commutation avec le cycle détermine toutes les colonnes de \(H\) à partir de l'une d'entre elles ; les douze puissances \(I,S,\ldots,S^{11}\) sont indépendantes et engendrent cet espace d'opérateurs. Par conséquent, \(y=O_Kh\), ce qui donne la seconde formule. Les deux premières bornes utilisent \(\|O_K^{-1}\|=1/\sqrt{589609}\). La dernière utilise \(\|O_{\delta y}\|\leq\|O_{\delta y}\|_{\rm F}
=\sqrt{12}\|\delta y\|\).
\end{proof}

Une réponse signifie ici un vecteur à douze composantes, non une
mesure scalaire. Par exemple, si \(H=2I+3S-S^5\), la réponse
\(y=HK\) restitue exactement les coefficients \(2,3,-1\) à leurs
positions au moyen de \eqref{md:antmem:k:filtro} ; l’inverseur a besoin de \(y,K,S\),
non des coefficients du filtre. En revanche, l’opérateur non nul dont la première
ligne est \((543,-234,0,\ldots,0)\) et dont les autres lignes sont nulles
annule \(K\). Une seule réponse n’identifie pas un opérateur arbitraire.

La fonction de référence admet une généralisation exacte. Dans un cycle de
\(n\) positions, \(H\mapsto Hg\) identifie tous les opérateurs
qui commutent avec le cycle si et seulement si
\(O_g=(g,Sg,\ldots,S^{n-1}g)\) est inversible. La suffisance est la
preuve précédente. Si \(O_g\) est singulier, un vecteur non nul de son noyau
donne un polynôme de degré inférieur à \(n\) qui annule \(g\) ; l’opérateur
est non nul car les puissances du cycle sont indépendantes. L’impulsion
unitaire a également une orbite orthonormale et sert de référence. La
conséquence spécifique pour \(K\) est que le même objet déjà utilisé dans
la clôture et l’incidence vérifie cette condition, sans avoir été choisi de
nouveau pour optimiser l’essai. On n’en déduit pas une identification de
dynamiques non linéaires ou d’états cachés arbitraires ; une application à
une linéarisation exige que celle-ci ait été définie et justifiée.

\begin{corollary}[Récupération de l'holonomie relative à partir du complément]
\label{md:antmem:k:holonomia}
Soient \(H_B\) connu et inversible et \(H_C\) un autre transport du même type. Le complément
\[
 D_\gamma=\frac{\sqrt8}{9}(H_C-H_B)
\]
détermine
\[
 H_B^{-1}H_C=I+\frac9{\sqrt8}H_B^{-1}D_\gamma.
\]
Si \(D_\gamma\) commute avec \(S\), en particulier si les deux
transports le font, \(y=D_\gamma K\) suffit à récupérer
\begin{equation}
 H_B^{-1}H_C
 =I+\frac9{\sqrt8}H_B^{-1}O_yO_K^{-1}.
 \label{md:antmem:k:holonomia-una-respuesta}
\end{equation}
L'erreur dans cette reconstruction est bornée par
\[
 \frac9{\sqrt8}\|H_B^{-1}\|
 \sqrt{\frac{12}{589609}}\,\|\delta y\|.
\]
\end{corollary}
\begin{proof}
On isole \(H_C\) dans la définition du complément et l'on applique \cref{md:antmem:k:operadores}. La borne résulte de la sous-multiplicativité.
\end{proof}

Pour \(H_B=S^3\) et \(H_C=S^4\), la donnée rationnelle
\(y/\sqrt8=(S^4-S^3)K/9\) permet de reconstruire l’holonomie relative
\(S\) sans fournir \(H_C\) à l’inverseur. Si \(H_B\) est
unitaire, le facteur d’erreur est d’environ \(0.01435510\). Sans
commutation, on maintient les douze réponses générales, non une hypothèse
fictive de symétrie. La sélection d’une boucle physique et sa représentation
appartiennent toujours à leur dynamique génératrice ; ce protocole ne les
remplace pas.

\subsubsection{Conséquence de la référence marquée pour deux orientations}
\label{md:antmem:orientaciones}
La conséquence utilisée dans les notes du cahier conserve le tableau
précédent comme antécédent. Pour un opérateur réel \(H\) qui commute avec
\(S\), soient \(y_+=HK\) et \(y_-=H^{\mathsf T}K\).
Leurs douze intensités de Fourier coïncident, mais \(y_+=y_-\) si et seulement
si \(H=H^{\mathsf T}\). Chaque réponse vectorielle complète, avec
\(K\) et l’orientation de \(S\), récupère son opérateur.
\begin{proof}
L’opérateur \(H\) est circulant ; ses multiplicateurs de Fourier sont
\(\lambda_m\), et ceux de \(H^{\mathsf T}\) sont leurs conjugués.
Ainsi les deux réponses ont pour intensités
\(|\lambda_m|^2|\widehat K_m|^2\).
Si leurs vecteurs coïncident, \((H-H^{\mathsf T})K=0\).
La commutation implique que \(H-H^{\mathsf T}\) annule également \(S^jK\)
pour tout \(j\). Ces vecteurs forment une base, donc
\(H-H^{\mathsf T}=0\). La récupération est
\(H=O_{y_+}O_K^{-1}\), et de même pour \(H^{\mathsf T}\).
\end{proof}
En particulier, les décalages \(S\) et \(S^{-1}\) donnent des réponses
de même spectre de puissance, mais distinctes. Les autocorrélations
exactes du tableau précédent donnent l’identité déjà utilisée dans le cahier :
\begin{equation}
 \|SK-S^{-1}K\|^2
 =2\bigl(\|K\|^2-\langle K,S^2K\rangle\bigr)
 =2\,175\,744>0.
 \label{md:antmem:distancia-orientada}
\end{equation}
Ce nombre est une conséquence des autocorrélations, non une donnée
de départ supplémentaire. Le contrôle rationnel du cahier récupère à partir des deux
réponses les coefficients \(e_1\) et \(e_{11}\), respectivement, par
inversion de \(O_K\). Une intensité spectrale ou un bilan scalaire ne
remplace pas la réponse vectorielle.
Ce corollaire concerne le support cyclique à douze composantes :
il n’identifie pas \(S\) à la dynamique enrichie complète et ne transforme
pas à lui seul \(S\leftrightarrow S^{-1}\) en conjugaison physique de charge.
Pour \(H\) général, la transposition ne signifie pas non plus une inversion temporelle ;
dans le cas \(S\), elle coïncide bien avec \(S^{-1}\).

\subsection{Couplage réversible et transport non stationnaire}
\label{md:antmem:acoplamiento}
\subsubsection{Extension réversible du bilan entre deux formes}

Soient \(B,C:V\to V'\) des isomorphismes qui conservent une forme non dégénérée \(b\) entre deux réalisations :

\[
B^*b'B=C^*b'C=b.
\]

La forme peut être une métrique positive ou une forme alternée. Dans les espaces réels, \(*\) désigne la transposition par rapport aux coordonnées choisies ; l’égalité exprime le transport de la forme entre les réalisations spécifiées. Dans les espaces de Hilbert, les opérateurs et leurs inverses sont bornés.


Avec \(s=\sqrt8\) et \(Q_V=\frac13\left(\begin{smallmatrix}sI&-I\\I&sI\end{smallmatrix}\right)\), nous définissons

\begin{equation}
\boxed{
\mathcal U(B,C)=Q_{V'}^*\operatorname{diag}(B,C)Q_V
=\frac19\begin{pmatrix}8B+C&\sqrt8(C-B)\\\sqrt8(C-B)&B+8C\end{pmatrix}.}
\label{md:antmem:acop:eq:1}
\end{equation}

Nous écrivons \(T=(8B+C)/9\), \(D=\sqrt8(C-B)/9\) et \(R=(B+8C)/9\). La mise à jour complète est

\begin{equation}
\binom{x'}{m'}=\begin{pmatrix}T&D\\D&R\end{pmatrix}\binom{x}{m}.
\label{md:antmem:acop:eq:2}
\end{equation}

\begin{theorem}
L’application \eqref{md:antmem:acop:eq:1} est réversible et conserve la forme \(b\oplus b\) :

\begin{equation}
\mathcal U(B,C)^*(b'\oplus b')\mathcal U(B,C)=b\oplus b,
\quad\mathcal U(B,C)^{-1}=\mathcal U(B^{-1},C^{-1}).
\label{md:antmem:acop:eq:3}
\end{equation}
\end{theorem}

\begin{proof}
On a \(Q^*Q=I\) par \(8+1=9\). Ses blocs scalaires font que \(Q\) préserve \(b\oplus b\) pour toute forme \(b\). L’application diagonale conserve les formes par hypothèse. La composition prouve \eqref{md:antmem:acop:eq:3}, et l’inversion des trois facteurs donne l’inverse déclaré.
\end{proof}

Pour une mémoire entrante \(m=0\), on retrouve le bilan \(b=T^*b'T+D^*b'D\). La seconde colonne détermine la réponse à une mémoire entrante arbitraire. La conservation des normes positives et la conservation de l’aire orientée sont les spécialisations respectives de \eqref{md:antmem:acop:eq:3}.

\paragraph{Réalisation collective dans les neuf copies
}

Dans \(V^9\), soit \(W_E=\operatorname{diag}(I,\ldots,I,E)\), avec \(E\) en neuvième position. Soient \(A(a,b)=(a/\sqrt8,\ldots,a/\sqrt8,b)\) et \(L=AQ\). L’inclusion uniforme est \(Jx=(x,\ldots,x)/3\). Alors \(L(x,0)=Jx\) et

\begin{equation}
W_E L=L\mathcal U(I,E).
\label{md:antmem:acop:eq:4}
\end{equation}

\begin{proof}
\(W_EA=A\operatorname{diag}(I,E)\) ; multiplier par \(Q\) et utiliser \(AQ=L\) donne \eqref{md:antmem:acop:eq:4}. L’orthogonalité de la moyenne des huit premières positions avec leurs sept différences montre en outre que le complément reste fixe sous \(W_E\).
\end{proof}

Soit \(P_{\rm cyc}(v_0,\ldots,v_8)=(v_8,v_0,\ldots,v_7)\). Le pas cyclique original est \(S_E=P_{\rm cyc}W_E\). Ainsi \(S_EL=(P_{\rm cyc}L)U(I,E)\), avec des coordinatisations d’entrée et de sortie distinctes. La réduction collective conserve le motif correspondant : \(S_E^9=I_9\otimes E\), tandis que \(W_E^9=\operatorname{diag}(I_8,E^9)\). Lorsqu’un lecteur utilise les neuf étiquettes individuelles, celles-ci sont conservées dans leur représentation originale.

\subsubsection{Composition exacte et réentrée de mémoire
}

Pour \(B_1,C_1:V_0\to V_1\) et \(B_2,C_2:V_1\to V_2\),

\begin{equation}
\boxed{\mathcal U(B_2,C_2)\mathcal U(B_1,C_1)
=\mathcal U(B_2B_1,C_2C_1).}
\label{md:antmem:acop:eq:5}
\end{equation}

\begin{proof}
Dans \eqref{md:antmem:acop:eq:1}, \(Q_{V_1}Q_{V_1}^*\) s’annulent, et les deux blocs diagonaux se multiplient dans l’ordre indiqué. La récurrence étend la formule à tout nombre fini de transports, en conservant leur ordre.

\end{proof}

Le bloc de lecture contient l’identité

\begin{equation}
\boxed{T_2T_1+D_2D_1=\frac{8B_2B_1+C_2C_1}{9}.}
\label{md:antmem:acop:eq:6}
\end{equation}

À partir de \((x,0)\), la première étape produit \((T_1x,D_1x)\). La seconde lecture est \(T_2T_1x+D_2D_1x\). Le second terme est exactement la contribution de la mémoire réintroduite. En stockant \(D_1x\) hors du couplage suivant et en poursuivant seulement avec \(T_1x\), on obtient le procédé à registres séparés, dont la lecture est \(T_2T_1x\). L’égalité \eqref{md:antmem:acop:eq:6} précise la différence entre les deux procédés.

La réalisation \(U\) conserve l’état du couplage et les produits ordonnés \(B_N\cdots B_1\), \(C_N\cdots C_1\). La conservation de tous les préfixes inclut en outre le registre des émissions de HMT, qui maintient chaque facteur avec son ordre.

\paragraph{Équation covariante d’incrément
}

Pour le registre séparé \(x_{n+1}=T_nx_n\), \(m_n=D_nx_n\), on a


\begin{equation}
x_{n+1}-B_nx_n=m_n/\sqrt8.
\label{md:antmem:acop:eq:7}
\end{equation}

Avec \(G_0=I\) et \(G_{n+1}=B_nG_n\), le télescopage donne


\begin{equation}
x_0=G_N^{-1}x_N-\frac1{\sqrt8}\sum_{n<N}G_{n+1}^{-1}m_n.
\label{md:antmem:acop:eq:8}
\end{equation}

La mémoire est l’incrément par rapport au transport entre formes \(B_n\), avec la normalisation déclarée. L’égalité \eqref{md:antmem:acop:eq:8} est une reconstruction finie algébrique ; son passage à une série infinie requiert la convergence. Le théorème suivant construit l’inverse stable au moyen de la limite d’opérateurs positifs.


\subsubsection{Mémoire d’une suite non stationnaire}

Dans une réalisation de Hilbert \(H\), soient \(C_n\) des unitaires, avec leur ordre de composition déclaré. Nous écrivons

\[
T_n=(8I+C_n)/9,\quad D_n=\sqrt8(C_n-I)/9,
\quad P_0=I,\quad P_{n+1}=T_nP_n,\quad A_N=P_N^*P_N.
\]

L’indice \(N\) compte les compressions avec stockage séparé, à la différence du transport complet avec réentrée. Pour des métriques variables avec \(B_n\) des isomorphismes isométriques, la conjugaison par \(G_n\) transporte les transports vers un espace de Hilbert fixe et produit ce même cas, avec \(\widetilde C_n=G_{n+1}^{-1}C_nG_n\). L’assertion utilise les métriques des fibres et les isométries déclarées entre elles.

\begin{theorem}[Terminal positif et récupération non stationnaire]
\label{md:antmem:terminal-no-estacionario}
Il existe un opérateur positif \(0\leq A_\infty\leq I\), limite forte de \(A_N\), et

\begin{equation}
\boxed{I=A_\infty+\sum_{n\ge0}P_n^*D_n^*D_nP_n.}
\label{md:antmem:acop:eq:9}
\end{equation}

L'application

\begin{equation}
\boxed{\mathcal Vv=(A_\infty^{1/2}v,(D_nP_nv)_{n\ge0})}
\label{md:antmem:acop:eq:10}
\end{equation}

est isométrique et possède un inverse à gauche stable


\begin{equation}
\mathcal V^*(a,(m_n))=A_\infty^{1/2}a+
\sum_{n\ge0}P_n^*D_n^*m_n,\qquad\|\mathcal V^*\|=1.
\label{md:antmem:acop:eq:11}
\end{equation}

\end{theorem}
\begin{proof}
L’identité locale \(T_n^*T_n+D_n^*D_n=I\) donne
\(A_N-A_{N+1}=P_N^*D_N^*D_NP_N\geq0\). Les formes quadratiques de
la suite positive décroissante convergent vers celle d’un opérateur positif
\(A_\infty\). Pour \(B_N=A_N-A_\infty\), l’inégalité
\(0\leq B_N^2\leq B_N\) transforme cette convergence en convergence forte.
Le télescopage fini et la limite prouvent \eqref{md:antmem:acop:eq:9}. L’évaluation sur \(v\)
prouve \eqref{md:antmem:acop:eq:10}. Son adjoint est \eqref{md:antmem:acop:eq:11} ; la série converge en norme parce que
les troncatures d’un registre de carré sommable convergent et que l’adjoint est borné.
\end{proof}

La reconstruction avec le terminal limite et \(N\) registres a pour erreur
exacte \((A_N-A_\infty)v\). Le vecteur \(A_\infty^{1/2}v\) est la
réalisation positive canonique du terminal de norme. Son existence provient
de la limite de Gram, indépendamment de la convergence vectorielle de \(P_Nv\).


\subsubsection{Trois propriétés différentes : épuisement, distinction et stabilité}

Soit \(Mv=(D_nP_nv)_n\). De \eqref{md:antmem:acop:eq:9}, \(M^*M=I-A_\infty\). Par conséquent :

\begin{enumerate}
\item Toute la norme de \(v\) passe au registre exactement lorsque \(A_\infty v=0\).

\item Le registre distingue chaque état exactement lorsque \(\ker(I-A_\infty)=\{0\}\).
\item La récupération à partir de la mémoire seule possède un inverse borné exactement lorsque
\(I-A_\infty\geq\varepsilon I\) pour un certain \(\varepsilon>0\).
Dans ce cas, un inverse à gauche est \((I-A_\infty)^{-1}M^*\).

\end{enumerate}

De plus,

\begin{equation}
\boxed{\ker M=\bigcap_{n\ge0}\operatorname{Fix}(C_n).}
\label{md:antmem:acop:eq:12}
\end{equation}

\begin{proof}
Si \(Mv=0\), alors \(D_0v=0\), d’où \(C_0v=v\)
et \(P_1v=v\). Par récurrence, \(D_nP_nv=0\) donne \(C_nv=v\) et
\(P_{n+1}v=v\). La réciproque est immédiate. Les trois équivalences se
déduisent du Gram \(M^*M\) et de la caractérisation des opérateurs
bornés inférieurement.
\end{proof}

\paragraph{Exemple exact qui sépare les trois propriétés}

Dans l’exemple algébrique \(C_0=-I\) et \(C_n=I\) pour \(n\geq1\),

\begin{equation}
A_\infty=\frac{49}{81}I,\qquad M^*M=\frac{32}{81}I,
\qquad m_0=-\frac{2\sqrt8}{9}v,
\qquad v=-\frac9{2\sqrt8}m_0.
\label{md:antmem:acop:eq:13}
\end{equation}

La mémoire contient \(32/81\) du carré de la norme et permet de récupérer le vecteur complet. Le terminal retient \(49/81\) et est un opérateur positif distinct d’un projecteur. La récupérabilité dépend du noyau et du conditionnement du lecteur. Lecture et mémoire constituent des composantes corrélées d’une isométrie.


Cet exemple sépare les trois conséquences logiques du théorème. L’application
à une trajectoire physique conserve en outre la sélection et la
provenance effective de ses \(C_n\) au moyen du TPK.

La lettre \(\eta\) de la réalisation gaussienne suivante désigne la
déformation de l’ellipse, c’est-à-dire le paramètre \(\zeta\) des
développements de transport ; ce n’est pas le coefficient d’action
\(\eta_{\rm ret}\). On maintient une section positive \(\hbar_a\) et
une même unité pour les deux coordonnées du quotient angulaire.
Les demi-axes conservés sont
\[
 a_S=\sqrt{2\hbar_a}\,e^{\eta/2},\qquad
 b_S=\sqrt{2\hbar_a}\,e^{-\eta/2},\qquad
 \eta=\operatorname{artanh}(C^*/A).
\]

\subsection{Covariance canonique et réciprocité de l'ellipse d'action}
\label{md:antmem:covarianza-accion}

La construction précédente détermine deux invariants complémentaires :
le produit des demi-axes fixe l’aire d’action et leur quotient conserve
l’anisotropie de la paire angulaire. Précisons maintenant la réalisation canonique
dans laquelle ces demi-axes sont deux écarts-types. La structure de
départ est l’ellipse déjà obtenue à partir des représentations angulaires et
d’action d’APP--TRIT--TPK ; la réalisation opératorielle s’établit après
cette construction.

Dans ce paragraphe, on écrit \(\eta=\eta_{\rm el}
=\operatorname{artanh}(C^*/A)\), avec \(|C^*/A|<1\) et
\(\hbar_a>0\). Le quotient angulaire utilise la même unité pour ses deux
composantes. Dans les coordonnées équilibrées \((Q,P)\), toutes deux de
dimension racine carrée d'action, définissons
\begin{equation}
 V_\eta:=\frac14
 \begin{pmatrix}a_S^2&0\\0&b_S^2\end{pmatrix}
 =\frac{\hbar_a}{2}
 \begin{pmatrix}e^\eta&0\\0&e^{-\eta}\end{pmatrix}.
 \label{md:antmem:covarianza-geometrica}
\end{equation}
Cette matrice positive possède des entrées de dimension action. Son interprétation
comme covariance quantique est spécifiée par la réalisation suivante.

\begin{proposition}[Réalisation gaussienne de la covariance d'action]
\label{md:antmem:covarianza-gaussiana}
Considérons la représentation canonique dans \(L^2(\mathbb R,dQ)\), avec
\(\widehat Qf(Q)=Qf(Q)\) et
\(\widehat Pf(Q)=-i\hbar_a f'(Q)\), initialement sur l'espace de
Schwartz \(\mathcal S(\mathbb R)\). L'état normalisé
\begin{equation}
 \psi_\eta(Q)=\bigl(\pi\hbar_a e^\eta\bigr)^{-1/4}
 \exp\!\left(-\frac{Q^2}{2\hbar_a e^\eta}\right)
 \label{md:antmem:gaussiana-accion}
\end{equation}
a des moyennes nulles et pour matrice de covariance \(V_\eta\). En particulier,
\begin{equation}
 (\Delta Q)^2=\frac{\hbar_a e^\eta}{2},\qquad
 (\Delta P)^2=\frac{\hbar_a e^{-\eta}}{2},\qquad
 \operatorname{Cov}(Q,P)=0,\qquad
 \det V_\eta=\frac{\hbar_a^2}{4}.
 \label{md:antmem:varianzas-accion}
\end{equation}
Cette réalisation sature l'inégalité de Robertson--Schrödinger.
\end{proposition}

\begin{proof}
Les opérateurs précédents sont essentiellement autoadjoints sur
\(\mathcal S(\mathbb R)\), qui constitue un cœur commun invariant ;
sur cet espace, on a
\([\widehat Q,\widehat P]=i\hbar_a I\). Pour
\(\widehat Y=(\widehat Q,\widehat P)\), la covariance symétrique est
\[
 V_{jk}=\frac12\left\langle
 (\widehat Y_j-\langle\widehat Y_j\rangle)
 (\widehat Y_k-\langle\widehat Y_k\rangle)
 +(\widehat Y_k-\langle\widehat Y_k\rangle)
 (\widehat Y_j-\langle\widehat Y_j\rangle)
 \right\rangle.
\]
Avec \(s=\hbar_a e^\eta>0\), soit
\(I_s=\int_{\mathbb R}e^{-Q^2/s}\,dQ\). Le produit de deux
intégrales identiques, évalué en coordonnées polaires, donne
\(I_s^2=2\pi\int_0^\infty r e^{-r^2/s}\,dr=\pi s\).
En outre, l'intégrale de
\(\frac{d}{dQ}(Qe^{-Q^2/s})\) est nulle, par décroissance aux deux
extrémités ; d'où
\(\int_{\mathbb R}Q^2e^{-Q^2/s}\,dQ=sI_s/2\).
Ces identités et la parité démontrent la normalisation, la moyenne nulle et
\(\langle\widehat Q^2\rangle=s/2\). De plus,
\(\widehat P\psi_\eta=i e^{-\eta}Q\psi_\eta\).
Ainsi, \(\langle\widehat P\rangle=0\),
\(\|\widehat P\psi_\eta\|^2=\hbar_a e^{-\eta}/2\) et
\(\operatorname{Re}\langle\widehat Q\psi_\eta,
\widehat P\psi_\eta\rangle=0\).

Pour tout état normalisé de ce domaine, Cauchy--Schwarz appliquée
à \((\widehat Q-\langle\widehat Q\rangle)\psi\) et
\((\widehat P-\langle\widehat P\rangle)\psi\) donne
\[
 (\Delta Q)^2(\Delta P)^2
 \geq \operatorname{Cov}(Q,P)^2
       +\frac14\left|\langle[\widehat Q,\widehat P]\rangle\right|^2
 =\operatorname{Cov}(Q,P)^2+\frac{\hbar_a^2}{4}.
\]
Les variances calculées atteignent l'égalité. De manière équivalente,
l'opérateur
\[
 a_\eta=\frac{e^{-\eta/2}\widehat Q
                  +i e^{\eta/2}\widehat P}{\sqrt{2\hbar_a}}
\]
satisfait \([a_\eta,a_\eta^\dagger]=I\) et
\(a_\eta\psi_\eta=0\), par substitution directe.
\end{proof}

\subsection{Covariance, forme normale et spectre gaussien}
\label{md:antmem:gauss:gaussianas-espectro}

La section positive d’action \(\hbar_a\) et la réalisation canonique
précédente déterminent la normalisation employée ici. On démontre la
correspondance entre covariance, spectre, pureté et entropie pour un
nombre fini de modes ; l’espace de Fock conserve toutes ses occupations.
La réduction symplectique et les formules gaussiennes sont des résultats
établis qui sont reproduits pour rendre autonome leur application au
transport de mémoire.

\subsubsection{Coordonnées canoniques et fonction caractéristique}
Dans \(L^2(\mathbb R^d)\), soient \(\widehat q_j=\widehat Q_j/\sqrt{\hbar_a}\),
\(\widehat p_j=\widehat P_j/\sqrt{\hbar_a}\) et
\(\widehat z=(\widehat q_1,\widehat p_1,\ldots,\widehat q_d,\widehat p_d)\).
Sur Schwartz,

\[
 [\widehat z_j,\widehat z_k]=iJ_{jk}I,\qquad
 J=\bigoplus_{j=1}^d\begin{pmatrix}0&1\\-1&0\end{pmatrix}.
\]
Cette convention fixe le signe de la forme canonique. Les opérateurs de
Weyl \(\mathcal W(\xi)=\exp(i\xi^{\mathsf T}\widehat z)\) sont définis
par les fermetures auto-adjointes des combinaisons linéaires indiquées.
Pour un état centré \(\rho\) de moments d’ordre deux finis, la
covariance normalisée et sa fonction caractéristique sont
\[
 W_{jk}=\operatorname{Tr}\rho(\widehat z_j\widehat z_k+\widehat z_k\widehat z_j),
 \qquad \chi_\rho(\xi)=\operatorname{Tr}(\rho\mathcal W(\xi)),\qquad
 V=\frac{\hbar_a}{2}W.
\]
L’état est gaussien centré lorsque

\begin{equation}
 \chi_\rho(\xi)=\exp(-\xi^{\mathsf T}W\xi/4).
 \label{md:antmem:gauss:eq:caracteristica-gaussiana}
\end{equation}
L’espérance de \((c^{\mathsf T}\widehat z)^\dagger
(c^{\mathsf T}\widehat z)\) est \(c^\dagger(W+iJ)c/2\) ; ainsi,
\(W+iJ\geq0\). De plus, \(W>0\) : si \(v\) est réel et
\(v^{\mathsf T}Wv=0\), la positivité implique \((W+iJ)v=0\),
donc \(Jv=0\) et \(v=0\).


\begin{lemma}[Unicité de la fonction caractéristique
]
\label{md:antmem:gauss:lem:unicidad-weyl}
Deux opérateurs de classe trace ayant la même fonction caractéristique sont égaux.
\end{lemma}
\begin{proof}
En regroupant les positions et les moments sous la forme \(\xi=(u,v)\),

\[
 (\mathcal W(u,v)f)(x)=e^{iu\cdot(x+v/2)}f(x+v).
\]
Pour un opérateur de rang fini à vecteurs de Schwartz et de noyau \(K_A\),

\[
 \operatorname{Tr}(A\mathcal W(u,v))=\int K_A(x+v/2,x-v/2)e^{iu\cdot x}\,dx.
\]
Plancherel et le changement de variables de jacobien absolu un donnent
\[
 \|A\|_{\rm HS}^2=(2\pi)^{-d}\int_{\mathbb R^{2d}}
 |\operatorname{Tr}(A\mathcal W(\xi))|^2\,d\xi.
\]
Les opérateurs précédents sont denses dans la classe trace. La convergence
dans cette norme implique la convergence en Hilbert--Schmidt et la convergence
uniforme des fonctions caractéristiques, car leur différence est
bornée par \(\|A-B\|_1\). L’égalité s’étend dans \(L^2\).
L’appliquer à la différence des opérateurs prouve l’assertion.
\end{proof}

\subsubsection{Réduction symplectique de la forme positive
}
\begin{theorem}[Forme normale de Williamson]
\label{md:antmem:gauss:thm:williamson}
Pour une matrice réelle symétrique \(W>0\) d’ordre \(2d\), il existe
\(S\in\operatorname{Sp}(2d,\mathbb R)\) et \(\nu_j>0\) tels que
\[
 W=SDS^{\mathsf T},\qquad D=\bigoplus_j\nu_jI_2.
\]
Les valeurs propres de \(JW\) sont \(\pm i\nu_j\). La condition
\(W+iJ\geq0\) équivaut à \(\nu_j\geq1\) pour tout \(j\).

\end{theorem}
\begin{proof}
La matrice \(B=W^{1/2}JW^{1/2}\) est antisymétrique et inversible.
Si \(-B^2u=\nu^2u\) et \(\|u\|=1\), le vecteur
\(v=-Bu/\nu\) est unitaire et orthogonal à \(u\) ; il satisfait
\(Bu=-\nu v\), \(Bv=\nu u\). Le complément de cette paire est
invariant. L’itération produit une matrice orthogonale \(O\) telle que
\(O^{\mathsf T}BO=JD\). En définissant \(S=W^{1/2}OD^{-1/2}\),

\[
 S^{\mathsf T}JS=J,\qquad SDS^{\mathsf T}=W.
\]
La similitude \(W^{1/2}(JW)W^{-1/2}=B\) détermine les valeurs
\(\nu_j\), avec multiplicité. La congruence par \(S^{-1}\)
transforme \(W+iJ\) en \(D+iJ\) ; chaque bloc a pour valeurs propres
\(\nu_j+1\) et \(\nu_j-1\).
\end{proof}

\begin{lemma}[Implémentation unitaire dans un nombre fini de modes]
\label{md:antmem:gauss:lem:implementacion-simplectica}
Toute matrice symplectique \(S\) admet un unitaire \(U_S\) qui satisfait
\(U_S^\dagger\mathcal W(\xi)U_S=\mathcal W(S^{\mathsf T}\xi)\).
La conjugaison \(\rho\mapsto U_S\rho U_S^\dagger\) transporte
la covariance par \(W\mapsto SWS^{\mathsf T}\).
\end{lemma}
\begin{proof}
Dans la décomposition polaire \(S=OP\), l’identité
\(J^{-1}(S^{\mathsf T}S)J=(S^{\mathsf T}S)^{-1}\), appliquée par
calcul spectral à la racine positive, prouve \(PJP=J\).
Alors \(O\) est orthogonale symplectique. Les espaces propres de \(P\)
s’apparient au moyen de \(J\) avec les valeurs propres \(\lambda,\lambda^{-1}\) ;
celui de valeur propre un est \(J\)-invariant. Une matrice
orthogonale symplectique \(K\) réduit donc \(P\) à
\(\bigoplus_j\operatorname{diag}(e^{r_j},e^{-r_j})\).

La dilatation est mise en œuvre par

\[
 (\mathcal D_rf)(x)=e^{-\frac12\sum_jr_j}
 f(e^{-r_1}x_1,\ldots,e^{-r_d}x_d).
\]
Un changement de variables prouve son unitarité et la conjugaison de
\((\widehat q_j,\widehat p_j)\) en
\((e^{r_j}\widehat q_j,e^{-r_j}\widehat p_j)\).
Une transformation orthogonale symplectique commute avec \(J\) et
correspond à une matrice \(u\in U(d)\) sur
\(a_j=(\widehat q_j+i\widehat p_j)/\sqrt2\). Sa seconde
quantification \(\Gamma(u)=\bigoplus_{n\geq0}u^{\otimes_s n}\)
est unitaire dans \(\mathcal F_s(\mathbb C^d)\), et l’identité sur les
produits symétriques finis donne la transformation des \(a_j\).

La base de Hermite identifie cet espace à \(L^2(\mathbb R^d)\).
Pour justifier sa complétude, si \(f\) est orthogonal à tous les
polynômes multipliés par \(e^{-|x|^2/2}\), la fonction entière
\(F(z)=\int f(x)e^{-|x|^2/2}e^{z\cdot x}\,dx\) a toutes
ses dérivées nulles en zéro. Ainsi, \(F=0\) ; en restreignant à
\(z=i\xi\), l’unicité de Fourier donne \(f=0\). Les
normalisations et l’orthogonalité résultent de la création et de l’annihilation.
La composition des mises en œuvre précédentes produit \(U_S\).
L’égalité de Weyl implique
\(\chi_{U_S\rho U_S^\dagger}(\xi)=\chi_\rho(S^{\mathsf T}\xi)\),
qui démontre la loi de covariance.

\end{proof}

\subsubsection{Calcul du spectre de la densité isotrope}
Dans un mode, les états cohérents normalisés sont
\[
 |z\rangle=e^{-|z|^2/2}\sum_{n\geq0}\frac{z^n}{\sqrt{n!}}|n\rangle.
\]
Leur fonction de position,
\(\pi^{-1/4}\exp[-(x-q_0)^2/2+ip_0(x-q_0/2)]\), où
\(q_0=\sqrt2\Re z\), \(p_0=\sqrt2\Im z\), donne par intégration
\[
 \langle z|\mathcal W(u,v)|z\rangle
 =e^{-(u^2+v^2)/4}e^{i(uq_0+vp_0)}.
\]
Pour \(s>0\), l’intégrale
\[
 \rho^{(s)}=\frac1{\pi s}\int_{\mathbb C} e^{-|z|^2/s}
 |z\rangle\langle z|\,d^2z
\]
converge en classe trace, est positive et a pour trace un. L’intégration
angulaire annule ses éléments hors de la diagonale d’occupation ; l’intégration
radiale donne
\[
 \langle n|\rho^{(s)}|n\rangle=\frac{s^n}{(1+s)^{n+1}},\qquad
 \chi_{\rho^{(s)}}(u,v)=e^{-(2s+1)(u^2+v^2)/4}.
\]
Avec \(\nu=2s+1\), l’unique état gaussien de covariance
\(\nu I_2\) est, par le lemme~\ref{md:antmem:gauss:lem:unicidad-weyl},
\begin{equation}
 \rho_\nu=(1-t)\sum_{n\geq0}t^n|n\rangle\langle n|,
 \qquad t=\frac{\nu-1}{\nu+1}.
 \label{md:antmem:gauss:eq:estado-geometrico}
\end{equation}
Le cas \(\nu=1\) est le projecteur sur le vide.

\begin{theorem}[Spectre, pureté et entropie
]
\label{md:antmem:gauss:thm:espectro-gaussiano}
Un état gaussien centré de covariance \(W\) et de valeurs
symplectiques \(\nu_1,\ldots,\nu_d\) est unitairement équivalent
à \(\bigotimes_j\rho_{\nu_j}\). Ses valeurs propres sont
\[
 \lambda_{\mathbf n}=\prod_j(1-t_j)t_j^{n_j},\qquad
 t_j=\frac{\nu_j-1}{\nu_j+1},\quad \mathbf n\in\mathbb N^d.
\]
Avec \(0\log0=0\),
\begin{align}
 \operatorname{Tr}\rho^2&=\prod_j\nu_j^{-1}=(\det W)^{-1/2},
 \label{md:antmem:gauss:eq:pureza-general}\\
 -\operatorname{Tr}(\rho\log\rho)&=\sum_j
 \left[\frac{\nu_j+1}{2}\log\frac{\nu_j+1}{2}
 -\frac{\nu_j-1}{2}\log\frac{\nu_j-1}{2}\right].
 \label{md:antmem:gauss:eq:entropia-general}
\end{align}
\end{theorem}
\begin{proof}
Williamson donne \(W=SDS^{\mathsf T}\). Le produit des densités
isotropes a pour fonction caractéristique \(e^{-\xi^{\mathsf T}D\xi/4}\).
L’implémentation unitaire et l’unicité de Weyl prouvent
l’équivalence et la formule spectrale. Dans un mode,
\[
 \sum_{n\geq0}(1-t)^2t^{2n}=\frac{1-t}{1+t}=\nu^{-1},
\]
et
\[
 -\sum_{n\geq0}(1-t)t^n\log((1-t)t^n)
 =-\log(1-t)-\frac{t}{1-t}\log t.
\]
La substitution \(t=(\nu-1)/(\nu+1)\) prouve les formules ; la
limite en \(t=0\) est nulle. Dans le produit fini, Tonelli justifie
la somme des termes d’entropie positifs ou nuls, et
\(\det W=\prod_j\nu_j^2\) donne l’expression déterminantale.

\end{proof}

Si l’état est la réduction d’un vecteur pur bipartite, ses
valeurs propres sont les carrés des coefficients de Schmidt.
La décomposition en valeurs singulières de l’opérateur de Hilbert--Schmidt défini
par ce vecteur démontre l’égalité des spectres non nuls des deux
réductions. L’entropie calculée est alors celle d’intrication.

\subsubsection{Préparation et transport des deux covariances}
L’application ultérieure utilise deux entrées gaussiennes pures identiques.
En général, pour \(M\in\operatorname{Sp}(2d,\mathbb R)\), leur covariance est
\(V_0=(\hbar_a/2)MM^{\mathsf T}\), et l’entrée conjointe possède la
covariance \(V_0\oplus V_0\). Le transport physique s’obtient
en conjuguant \(U_H\) par \(M\oplus M\). Dans la coordinatisation équilibrée,
la covariance normalisée d’entrée est \(I_{4d}\) ; avec
\[
 T=\frac{8I+H}{9},\qquad D=\frac{\sqrt8(H-I)}9,
\]
la première réduction a
\[
 W=TT^{\mathsf T}+DD^{\mathsf T}=\frac{8I+HH^{\mathsf T}}9.
\]
Sa fonction caractéristique s’obtient en annulant les arguments du
second groupe de modes. Elle reste donc gaussienne et les théorèmes
précédents calculent tout son spectre, sans tronquer le registre d’occupation.

\subsection{Prolongement du transport de mémoire : réalisation quantique}

La réalisation quantique utilise le couplage réversible de la section~\ref{md:antmem:acoplamiento}. Nous considérons d’abord un plan canonique réel de matrice alternée \(\Omega=\left(\begin{smallmatrix}0&1\\-1&0\end{smallmatrix}\right)\). La partition nonadique détermine

\begin{equation}
Q=\frac13\begin{pmatrix}\sqrt8 I&-I\\I&\sqrt8 I\end{pmatrix},\qquad
U_H=Q^{\mathsf T}\operatorname{diag}(I,H)Q
=\frac19\begin{pmatrix}8I+H&\sqrt8(H-I)\\
\sqrt8(H-I)&I+8H\end{pmatrix}.
\label{md:antmem:cuantica:eq:5.1}
\end{equation}

L’opérateur \(H\in\operatorname{Sp}(2,\mathbb R)\) transporte ce plan. La matrice \(Q\) conserve le produit euclidien et la forme \(\Omega\oplus\Omega\) ; c’est pourquoi \(U_H\) est symplectique. La première composante définit le système observé et la seconde le registre de mémoire ; la transformation complète agit sur les deux.

Le lacet ordonné dont les identités sont reproduites dans
\ref{md:antmem:lazo-ep} est

\begin{equation}
C=\begin{pmatrix}0&-1\\1&-1\end{pmatrix},\quad
P=\begin{pmatrix}1&1\\0&1\end{pmatrix},\quad
H_n=C^{-1}P^{-n}CP^n
=\begin{pmatrix}1+n&n^2\\n&n^2-n+1\end{pmatrix}.
\label{md:antmem:cuantica:eq:5.2}
\end{equation}

Pour chaque \(n\in\mathbb Z\), on a \(\det H_n=1\). En particulier, pour \(n=1\),

\begin{equation}
H_1=\begin{pmatrix}2&1\\1&1\end{pmatrix}
=F_1^2,\quad F_1=\begin{pmatrix}1&1\\1&0\end{pmatrix}
=P F_{\rm av}P^{-1}.
\label{md:antmem:cuantica:eq:5.3}
\end{equation}

Son spectre est \(\{\varphi^2,\varphi^{-2}\}\), par la reconnaissance ultérieure du mode d’or déjà engendré. L’entier \(n\) spécifie le nombre orienté d’itérations dans le lacet choisi. La collection des transports est déterminée algébriquement ; sa réalisation comme protocole physique conserve en outre les données de préparation et d’évolution correspondantes.

\subsubsection{Identités algébriques du lacet et reconnaissance ultérieure}
\label{md:antmem:lazo-ep}
Pour conserver la preuve des identités utilisées, on reproduit
leur composition matricielle. Dans la même coordinatisation,
\[
 N_0=\begin{pmatrix}0&1\\0&0\end{pmatrix}.
\]
Soit
\[
 P=I+N_0=\begin{pmatrix}1&1\\0&1\end{pmatrix}.
\]
Comme \(N_0^2=0\), on a \(P^n=I+nN_0\) pour tout
\(n\in\mathbb Z\).
Le lacet ordonné de quatre transports donne

\begin{equation}
H_n=C^{-1}P^{-n}CP^n
=\begin{pmatrix}1+n&n^2\\n&n^2-n+1\end{pmatrix}.
\label{md:antmem:matriz:eq:28}
\end{equation}

Pour \(n\ne0\), sa différence normalisée est

\[
F_n=\frac{H_n-I}{n}
=\begin{pmatrix}1&n\\1&n-1\end{pmatrix}.
\]

En multipliant ces matrices, on obtient, sans introduire de valeurs spectrales,

\begin{equation}
\boxed{F_n^2=H_n,\qquad F_n^2-nF_n-I=0.}
\label{md:antmem:matriz:eq:29}
\end{equation}

Avec \(H_B=I\), la mémoire est
\(D_{\gamma,n}=(\sqrt8/9)nF_n\).
Par conséquent,

\begin{equation}
\boxed{H_n=\left(\frac9{n\sqrt8}D_{\gamma,n}\right)^2.}
\label{md:antmem:matriz:eq:30}
\end{equation}

La réponse matricielle de mémoire reconstruit l’auto-échelle
hyperbolique. Pour \(n>0\), la reconnaissance spectrale ultérieure
identifie

\begin{equation}
\rho_{{\rm ep},n}=\frac{n+\sqrt{n^2+4}}2,\quad
\operatorname{Spec}F_n=\{\rho_{{\rm ep},n},-\rho_{{\rm ep},n}^{-1}\},\quad
\operatorname{Spec}H_n=\{\rho_{{\rm ep},n}^2,\rho_{{\rm ep},n}^{-2}\}.
\label{md:antmem:matriz:eq:31}
\end{equation}

Le cas \(n=1\) produit le polynôme d’or.
La relation avec l’opérateur
\(F_{\rm av}=
\begin{pmatrix}0&1\\1&1\end{pmatrix}\)
engendré par le calendrier conserve le changement
de coordinatisation :

\begin{equation}
F_1=PF_{\rm av}P^{-1},\qquad
H_1=PF_{\rm av}^{\,2}P^{-1}.
\label{md:antmem:matriz:eq:32}
\end{equation}

Le transport \(P\), de déterminant un, entrelace les deux
matrices et conserve \(\Omega\).

\subsubsection{État initial, transport et lecture}

La proposition~\ref{md:antmem:covarianza-gaussiana} construit la réalisation gaussienne de l’ellipse d’action. Dans la section positive \(\hbar_a>0\), avec \(A>0\) et \(|C^*/A|<1\), sa covariance est


\begin{equation}
V_\eta=\frac{\hbar_a}{2}M_\eta M_\eta^{\mathsf T},\qquad
M_\eta=\operatorname{diag}(e^{\eta/2},e^{-\eta/2}),\qquad
\eta=\operatorname{artanh}(C^*/A).
\label{md:antmem:cuantica:eq:5.4}
\end{equation}

Dans la représentation de Schrödinger sur \(L^2(\mathbb R,dQ)\), avec \(\widehat Q\) la multiplication par \(Q\) et \(\widehat P=-i\hbar_a\,d/dQ\), la relation \([\widehat Q,\widehat P]=i\hbar_a I\) vaut dans l’espace de Schwartz. L’état pur normalisé est

\[
\psi_\eta(Q)=(\pi\hbar_a e^\eta)^{-1/4}
\exp[-Q^2/(2\hbar_a e^\eta)].
\]

La préparation est le produit \(\psi_\eta\otimes\psi_\eta\). Le transport de \eqref{md:antmem:cuantica:eq:5.1} s’exprime dans la même coordinatisation de covariance par

\begin{equation}
\widetilde U=(M_\eta\oplus M_\eta)U_H
(M_\eta^{-1}\oplus M_\eta^{-1}).
\label{md:antmem:cuantica:eq:5.5}
\end{equation}

Le lemme~\ref{md:antmem:gauss:lem:implementacion-simplectica} fournit une implémentation unitaire de cette matrice symplectique à deux modes. Les implémentations qui diffèrent d’une phase globale produisent la même conjugaison de l’état. La covariance se transforme par \(V\mapsto\widetilde U V\widetilde U^{\mathsf T}\) ; l’état conjoint reste pur et sa trace partielle sur le second mode détermine la lecture visible, conformément à la construction de la section~\ref{md:antmem:gauss:gaussianas-espectro}.

\begin{theorem}[Mélange réduit par déformation relative]

La covariance de la lecture visible est

\begin{equation}
\boxed{V_{\rm vis}=
\frac{\hbar_a}{2}M_\eta
\frac{8I+HH^{\mathsf T}}9
M_\eta^{\mathsf T}.}
\label{md:antmem:cuantica:eq:5.6}
\end{equation}

Sa valeur symplectique normalisée est

\begin{equation}
\boxed{\nu(H):=\frac{2\sqrt{\det V_{\rm vis}}}{\hbar_a}
=\sqrt{1+\frac8{81}\left[\operatorname{tr}(HH^{\mathsf T})-2\right]}.}
\label{md:antmem:cuantica:eq:5.7}
\end{equation}
\end{theorem}

\begin{proof}
Soient \(T=(8I+H)/9\) et \(D=\sqrt8(H-I)/9\). Les deux modes initiaux sont indépendants et ont la même covariance. Dans la coordinatisation équilibrée de l’état initial, la covariance visible normalisée est \(TT^{\mathsf T}+DD^{\mathsf T}\). Le développement donne

\[
\begin{aligned}
TT^{\mathsf T}+DD^{\mathsf T}
&=\frac{64I+8(H+H^{\mathsf T})+HH^{\mathsf T}
+8[HH^{\mathsf T}-H-H^{\mathsf T}+I]}{81}\\
&=\frac{8I+HH^{\mathsf T}}9.
\end{aligned}
\]

Cela démontre \eqref{md:antmem:cuantica:eq:5.6}. Puisque \(\det M_\eta=1\) et \(\det(HH^{\mathsf T})=1\),
\(\det(8I+HH^{\mathsf T})=65+8\operatorname{tr}(HH^{\mathsf T})\).
La division par \(81\) démontre \eqref{md:antmem:cuantica:eq:5.7}. Les deux valeurs propres positives de \(HH^{\mathsf T}\) sont réciproques ; leur somme est au moins deux. En conséquence, \(\nu\geq1\), avec égalité exactement lorsque \(HH^{\mathsf T}=I\).
\end{proof}

L’échelle \(\hbar_a\) et la déformation \(\eta\) sont conservées jusqu’au calcul du déterminant. Leur simplification est covariante car l’état et l’opérateur sont transportés conjointement. L’expression \(HH^{\mathsf T}\) utilise le produit euclidien de la coordinatisation équilibrée de l’état initial ; un changement de coordinatisation transporte également ce produit métrique.

Un transport \(H\) orthogonal et symplectique peut avoir \(D\neq0\) tout en conservant le produit gaussien isotrope, avec \(\nu=1\). Le complément opératoire et l’intrication de l’état ont donc des critères distincts. Dans cette préparation, la déformation relative non isométrique caractérise exactement le cas \(\nu>1\).

\subsubsection{Évaluation exacte du lacet d’or}

Pour \eqref{md:antmem:cuantica:eq:5.2},

\begin{equation}
\operatorname{tr}(H_nH_n^{\mathsf T})-2
=n^2(2n^2-2n+5),\qquad
\nu_n^2=1+\frac8{81}n^2(2n^2-2n+5).
\label{md:antmem:cuantica:eq:5.8}
\end{equation}

L’identité résulte de la somme des carrés des quatre entrées. Pour \(n=1\), la somme est sept et

\begin{equation}
\boxed{\nu_1=\frac{11}{9},\quad
\operatorname{Tr}\rho_{\rm vis}^2=\frac9{11},\quad
\bar n_{\rm W}=\frac19,\quad
\frac{S_{\rm vis}}{k_B}=\frac{10}{9}\log10-\log9.}
\label{md:antmem:cuantica:eq:5.9}
\end{equation}

Le théorème~\ref{md:antmem:gauss:thm:espectro-gaussiano} démontre que l’état gaussien centré d’un mode de valeur symplectique \(\nu\) est unitairement équivalent à la densité géométrique de \eqref{md:antmem:gauss:eq:estado-geometrico}, dont l’occupation moyenne est \(\bar n_{\rm W}=(\nu-1)/2\). Pour \(\nu=11/9\), ses valeurs propres sont

\begin{equation}
\lambda_j=\frac9{10}\,10^{-j},\qquad j=0,1,2,\ldots.
\label{md:antmem:cuantica:eq:5.10}
\end{equation}

La série a pour somme un. Son carré a pour somme
\((9/10)^2/(1-10^{-2})=9/11\). Enfin,


\[
-\sum_j\lambda_j\log\lambda_j
=-\log(9/10)+\log10\sum_jj\lambda_j
=\frac{10}{9}\log10-\log9.
\]

La réduction gaussienne et son spectre sont démontrés dans la section~\ref{md:antmem:gauss:gaussianas-espectro} ; les deux séries précédentes évaluent exactement la pureté et l’entropie de cette réalisation. Comme l’état conjoint est pur, \(S_{\rm vis}/k_B\) est également son entropie adimensionnelle d’intrication. La transformation globale est unitaire, tandis que la restriction au premier mode a une entropie positive.

Le résultat correspond à la préparation et au transport spécifiés. Son identification à une mesure du vide requiert la réalisation expérimentale de ce protocole. L’entropie d’une réduction et la chaleur transférée restent des grandeurs distinctes.

